% ===== PRÉAMBULE CANONIQUE — à coller VERBATIM en tête de CHAQUE .tex =====
\documentclass[11pt,a4paper]{article}
\usepackage[utf8]{inputenc}\usepackage[T1]{fontenc}\usepackage{lmodern}
\usepackage[french]{babel}
\usepackage{amsmath,amssymb,mathtools}\usepackage{icomma}
\usepackage[version=4]{mhchem}          % \ce{...} : équations & formules
\usepackage{siunitx}\sisetup{locale=FR} % \qty{}{}, \si{}, \ang{}
\usepackage{chemfig}                    % \chemfig{...} : structures développées
\usepackage{xcolor}\usepackage{colortbl}
\usepackage{tikz}\usepackage{pgfplots}\pgfplotsset{compat=1.18}
\usepackage[most]{tcolorbox}\usepackage{enumitem}\usepackage{array,booktabs}
\usepackage[a4paper,margin=2.2cm]{geometry}
\usetikzlibrary{arrows.meta,calc,angles,quotes,positioning,patterns,backgrounds,shapes.geometric}
\definecolor{primary}{RGB}{222,49,99}\definecolor{primarydark}{RGB}{150,22,55}
\definecolor{defcol}{RGB}{0,114,178}\definecolor{methcol}{RGB}{52,92,148}
\definecolor{excol}{RGB}{0,135,90}\definecolor{attncol}{RGB}{200,40,55}
\tcbset{boxbase/.style={enhanced,breakable,boxrule=0.9pt,arc=2.5pt,
  left=3mm,right=3mm,top=2.3mm,bottom=2.3mm,fonttitle=\bfseries,coltitle=white}}
% --- boîtes TUTORIEL (noms EXACTS, ne pas renommer) ---
\newtcolorbox{cle}[1][]{boxbase,colback=defcol!6,colframe=defcol,
  title={\textsc{À retenir}\ifx\relax#1\relax\relax\else\textmd{ : \itshape #1}\fi}}
\newtcolorbox{methode}[1][]{boxbase,colback=methcol!7,colframe=methcol,sharp corners,
  title={\textsc{Méthode}\ifx\relax#1\relax\relax\else\textmd{ --- \itshape #1}\fi}}
\newtcolorbox{exemplebox}[1][]{boxbase,colback=excol!5,colframe=excol,
  title={\textsc{Exemple}\ifx\relax#1\relax\relax\else\textmd{ : \itshape #1}\fi}}
\newtcolorbox{piege}[1][]{boxbase,colback=attncol!6,colframe=attncol,
  title={\textsc{Piège}\ifx\relax#1\relax\relax\else\textmd{ : \itshape #1}\fi}}
\newtcolorbox{remarque}[1][]{boxbase,colback=black!4,colframe=black!45,
  title={\textsc{Remarque}\ifx\relax#1\relax\relax\else\textmd{ : \itshape #1}\fi}}
% --- boîtes EXERCICE / CORRIGÉ (noms EXACTS, auto-comptées) ---
\newtcolorbox[auto counter]{exo}[1][]{boxbase,colback=primary!4,colframe=primary,
  title={\textsc{Exercice}~\thetcbcounter\ifx\relax#1\relax\relax\else\textmd{ --- \itshape #1}\fi}}
\newtcolorbox[auto counter]{corrige}[1][]{boxbase,colback=excol!4,colframe=excol,
  title={\textsc{Corrigé}~\thetcbcounter\ifx\relax#1\relax\relax\else\textmd{ --- \itshape #1}\fi}}
\newcommand{\R}{\mathbb{R}}\newcommand{\N}{\mathbb{N}}\newcommand{\Z}{\mathbb{Z}}\newcommand{\ds}{\displaystyle}
% (un \usepackage{fancyhdr}+en-tête est possible ; il est ignoré côté web)
% =========================================================================
\begin{document}
\sisetup{per-mode=symbol}
\begin{corrige}[juste ou fidèle ?]
\emph{Fidèle} $=$ résultats serrés entre eux ; \emph{juste} $=$ centrés sur la vraie valeur ($10,00$).
\begin{enumerate}[label=\alph*)]
  \item Serrés et centrés sur $10,00$ : \textbf{juste \emph{et} fidèle}.
  \item Serrés mais autour de $10,50$ : \textbf{fidèle mais non juste} (biais de $+0,50$).
  \item Dispersés mais de moyenne $\approx 10,0$ : \textbf{juste (en moyenne) mais non fidèle}.
  \item Dispersés et décentrés : \textbf{ni juste ni fidèle}.
\end{enumerate}
\end{corrige}

\begin{corrige}[incertitude absolue sous-entendue]
$\Delta x$ vaut une unité du rang du dernier chiffre significatif.
\begin{enumerate}[label=\alph*)]
  \item $\Delta = \qty{0,01}{\gram}$.
  \item $\Delta = \qty{0,1}{\milli\litre}$.
  \item $\Delta = \qty{0,1}{\gram}$.
  \item $\Delta = \qty{0,001}{\second}$.
\end{enumerate}
\end{corrige}

\begin{corrige}[encadrer la vraie valeur]
On prend $[\,x-\Delta x\,;\ x+\Delta x\,]$.
\begin{enumerate}[label=\alph*)]
  \item $[\,\qty{13,01}{\gram}\,;\ \qty{13,03}{\gram}\,]$.
  \item $[\,\qty{5,55}{\centi\metre}\,;\ \qty{5,65}{\centi\metre}\,]$.
  \item $[\,\qty{20,95}{\milli\litre}\,;\ \qty{20,97}{\milli\litre}\,]$.
\end{enumerate}
\end{corrige}

\begin{corrige}[incertitude relative]
On divise $\Delta x$ par $x$, puis on convertit en pourcentage.
\begin{enumerate}[label=\alph*)]
  \item $\dfrac{0,01}{13,02}=7,7\times10^{-4}\approx \boxed{0,077\,\%}$.
  \item $\dfrac{0,1}{4,5}=2,2\times10^{-2}\approx \boxed{2,2\,\%}$.
  \item $\dfrac{0,1}{250,0}=4,0\times10^{-4}\approx \boxed{0,040\,\%}$.
\end{enumerate}
Même incertitude absolue ($\qty{0,1}{}$), mais la mesure (c) est bien plus précise en \emph{relatif} :
c'est ce que traduit son plus grand nombre de CS.
\end{corrige}

\end{document}
